% exercises.tex -- irregular verbs, word formation and phrasal verbs
%
%   make build      -> exercises.pdf (worksheet) and exercises-solutions.pdf
%
% Layout: a key-terms summary comes first. Each topic (irregular verbs, word
% formation, phrasal verbs) starts on a new page, and within a topic the
% exercises go from the easiest type (matching) to the hardest (writing the
% forms yourself).
%
\documentclass[11pt]{article}
\usepackage{worksheet}

\begin{document}

\worksheetheader[Irregular verbs, nouns and phrasal verbs]{English Worksheet}

% ===========================================================================
% Key terms
% ===========================================================================
\begin{keyterms}
  \keyterm{Infinitive}
    {the basic form of a verb, usually with \textit{to}:
     \textit{to sing}.}
  \keyterm{Past simple}
    {used for finished actions in the past:
     \textit{Yesterday I sang in the choir.}}
  \keyterm{Past participle}
    {used with \textit{have}, \textit{has} or \textit{had}:
     \textit{I have sung in the choir since 2020.}}
  \keyterm{Irregular verb}
    {a verb whose past forms do not end in \textit{-ed}.
     You have to learn them by heart.}

  \begin{formtable}{Infinitive, Past simple, Past participle}
    \formrow{to hit,   hit,   hit}
    \formrow{to sleep, slept, slept}
    \formrow{to sing,  sang,  sung}
  \end{formtable}

  \keyterm{Noun}
    {a word for a person, a thing or an activity. Many nouns are made
     from verbs: \textit{to teach} $\rightarrow$ \textit{the teacher},
     \textit{teaching}.}
  \keyterm{Phrasal verb}
    {a verb plus a small word. Together they have a new meaning:
     \textit{to look after} = to take care of.}
  \keyterm{Particle}
    {the small word in a phrasal verb, for example \textit{up}, \textit{on},
     \textit{off}, \textit{out}.}
\end{keyterms}

\clearpage

% ===========================================================================
% Topic 1: irregular verbs
% ===========================================================================

% ---------------------------------------------------------------------------
\exercise{Irregular verbs: matching}%
  {Draw a line from each infinitive to its past participle.}

\begin{matching}{Infinitive}{Past participle}
  \matchpair{to go}{gone}
  \matchpair{to catch}{caught}
  \matchpair{to swear}{sworn}
  \matchpair{to grow}{grown}
  \matchpair{to melt}{melted}
  \matchpair{to wear}{worn}
  \matchpair{to draw}{drawn}
\end{matching}

% ---------------------------------------------------------------------------
\exercise{Irregular verbs}{Complete the table.}

\begin{formtable}{Infinitive, Past simple, Past participle}
  \formrow{to go,        \ans{went},    \ans{gone}}
  \formrow{to catch,     \ans{caught},  \ans{caught}}
  \formrow{to swear,     \ans{swore},   \ans{sworn}}
  \formrow{\ans{to grow}, grew,         \ans{grown}}
  \formrow{to melt,      \ans{melted},  \ans{melted}}
  \formrow{\ans{to wear}, \ans{wore},   worn}
  \formrow{to draw,      \ans{drew},    \ans{drawn}}
\end{formtable}

% ---------------------------------------------------------------------------
\exercise{Fill in the gaps}{Put the verb in brackets into the correct form.}

\begin{gapexercise}
  \item Yesterday we \gapverb{to go}{went} to the cinema.
  \item Have you ever \gapverb{to catch}{caught} a fish with your bare hands?
  \item He \gapverb{to swear}{swore} that he would never tell anyone.
  \item Our tomatoes have \gapverb{to grow}{grown} very quickly this summer.
  \item By the time we got home, the ice cream had \gapverb{to melt}{melted}.
  \item She \gapverb{to wear}{wore} her new coat to the party last night.
  \item The children have \gapverb{to draw}{drawn} a map of the village.
  \item Where is Tom? He has \gapverb{to go}{gone} to the shops.
\end{gapexercise}

\clearpage

% ===========================================================================
% Topic 2: word formation
% ===========================================================================

% ---------------------------------------------------------------------------
\exercise{From verb to noun}{Write the person and the activity (or result).}

\begin{formtable}{Verb, Person, Activity / result}
  \formrow{to cycle,       the cyclist,     cycling}
  \formrow{to type,        \ans{the typist},  \ans{typing}}
  \formrow{to grow,        \ans{the grower},  \ans{growth}}
  \formrow{to catch,       \ans{the catcher}, \ans{the catch}}
  \formrow{\ans{to sail},  the sailor,      \ans{sailing}}
  \formrow{to visit,       \ans{the visitor}, \ans{the visit}}
  \formrow{\ans{to drive}, the driver,      \ans{driving}}
\end{formtable}

\clearpage

% ===========================================================================
% Topic 3: phrasal verbs
% ===========================================================================

% ---------------------------------------------------------------------------
\exercise{Phrasal verbs: matching}%
  {Draw a line from each phrasal verb to its meaning.}

\begin{matching}[match-sep=1.35cm]{Phrasal verb}{Meaning}
  \matchpair{to put on}{to dress yourself in clothes}
  \matchpair{to take off}{to remove clothes; (of a plane) to leave the ground}
  \matchpair{to go on}{to continue}
  \matchpair{to catch up with}{to reach someone who is ahead of you}
  \matchpair{to grow up}{to become an adult}
  \matchpair{to wear out}{to use something until it can't be used any more}
  \matchpair{to draw up}{to prepare a plan or a document}
  \matchpair{to swear by}{to trust something completely}
  \matchpair{to melt away}{to disappear slowly}
\end{matching}

% ---------------------------------------------------------------------------
\exercise{Phrasal verbs: particles}%
  {Fill in the missing word. You can use some words more than once.}

\wordbox{on, off, up, out, by, with, away}

\begin{gapexercise}
  \item It's cold outside -- put \gap[2.5cm]{on} your gloves.
  \item Our plane took \gap[2.5cm]{off} two hours late.
  \item My grandmother swears \gap[2.5cm]{by} ginger tea for a cold.
  \item These trainers are completely worn \gap[2.5cm]{out}.
  \item Hurry up! We need to catch up \gap[2.5cm]{with} the others.
  \item When I grow \gap[2.5cm]{up}, I want to be a pilot.
  \item The lawyer drew \gap[2.5cm]{up} a new contract.
  \item After the concert, the crowd slowly melted \gap[2.5cm]{away}.
  \item Please go \gap[2.5cm]{on} -- I'm listening.
\end{gapexercise}

% ---------------------------------------------------------------------------
\exercise{Your turn}%
  {Write two sentences of your own. Use \textit{to wear out} and \textit{to draw up}.}

\writinglines{4}

\end{document}
